\section{Introduction}

Genetic algorithms (GAs) are named for an evolutionary metaphor, but the
population size $N$ is almost always set by heuristics internal to the GA
literature -- schema-survival and noise-to-signal arguments
\cite{goldberg1991sizing} -- rather than by population genetics' own
quantitative account of when a population is ``too small'' for selection
to matter. That account is precise: in a finite haploid population, a new
mutation with selection advantage $s$ fixes with probability given by
Kimura's diffusion approximation \cite{kimura1962probability},
\[
u(s,N) = \frac{1-e^{-2s}}{1-e^{-2Ns}},
\]
which is close to the neutral value $1/N$ when $Ns \ll 1$ (drift
dominates) and close to $2s$ when $Ns \gg 1$ (selection dominates). The
crossover, $N \approx 1/s$, is a specific, testable prediction for a
critical population size $N_c$.

We ask two questions about a GA searching an NK fitness landscape
\cite{kauffman1987towards}, whose epistasis parameter $K$ tunes
ruggedness (the number of local optima) directly: (1) is there a
population size $N_c(K)$ below which the GA reliably fails to find the
global optimum, and above which it reliably succeeds; and (2) if so, does
$N_c(K)$ grow with $K$ for the reason population genetics would predict --
a shrinking single-mutation selection coefficient $s(K)$, via
$N_c \approx 1/s(K)$?

We find a clear answer to (1) and a clear, negative answer to (2): a
real threshold exists and grows sharply with $K$, but the mechanism
proposed for it is falsified by direct measurement.
